\appendix

\section{Formal Models and Pseudocode}

This appendix documents formal specifications and pseudocode for key server operations. Each subsection specifies the algorithm's purpose, formal semantics, and role in the overall design.

\subsection{Type-Checking Algorithm for White-Box Programs}

\textbf{Purpose}: Verify that a program respects information flow constraints before execution. Type-checker enforces explicit flow rules (label ordering) and implicit flow rules (program counter tracking through conditionals). If type-checking succeeds, program certified secure; if fails, program rejected and not executed.

\textbf{Formal Specification}:

Algorithm TYPECHECK($P$, $L_{\text{in}}$, $L_{\text{out}}$):

\begin{algorithmic}
\State Input: Program $P$, input label environment $L_{\text{in}}: \{\text{var} \to \text{label}\}$, expected output label $L_{\text{out}}$
\State Initialize: variable label environment $\Gamma = L_{\text{in}}$, program counter label $L_{pc} = \text{public}$
\State Initialize: errors = []

\Function{typecheck\_statement}{$s$}
  \If{$s$ is assignment $x := e$}
    \State $L_e \gets$ \Call{typecheck\_expr}{$e$}
    \State $L_x \gets \Gamma[x]$ \textit{\# Get label for variable $x$}
    \State $L_{\text{combined}} \gets L_{pc} \sqcup L_e$ \textit{\# Combine with program counter}
    \If{$\neg (L_{\text{combined}} \leq L_x)$}
      \State append("EXPLICIT\_FLOW\_VIOLATION: $L_{\text{combined}} \not\leq L_x$" to errors)
    \EndIf
    \State $\Gamma[x] \gets L_x$ \textit{\# Update label (unchanged if check passed)}
  \ElsIf{$s$ is conditional if($c$) $\{ S_{\text{true}} \}$ else $\{ S_{\text{false}} \}$}
    \State $L_c \gets$ \Call{typecheck\_expr}{$c$}
    \State $L_{pc}^{\text{old}} \gets L_{pc}$
    \State $L_{pc} \gets L_{pc}^{\text{old}} \sqcup L_c$ \textit{\# Update pc: join with condition label}
    \State \Call{typecheck\_statement}{$S_{\text{true}}$}
    \State \Call{typecheck\_statement}{$S_{\text{false}}$}
    \State $L_{pc} \gets L_{pc}^{\text{old}}$ \textit{\# Restore pc after branch}
  \ElsIf{$s$ is call declassify($x$, principal, $R'$)}
    \State $L_x \gets \Gamma[x]$ \textit{\# Current label of $x$}
    \State \textit{\# Extract owner from label $L_x = o: R$}
    \State $o \gets$ \Call{owner}{$L_x$}
    \If{principal $\neq o$}
      \State append("DECLASSIFY\_UNAUTHORIZED: principal $\neq$ owner" to errors)
    \Else
      \State \textit{\# Owner authorizes declassification}
      \State $R \gets$ \Call{readers}{$L_x$}
      \If{$\neg (R \subseteq R')$}
        \State append("DECLASSIFY\_SHRINKING: $R \not\subseteq R'$" to errors)
      \Else
        \State \textit{\# Approved: update label}
        \State $\Gamma[x] \gets o: R'$
      \EndIf
    \EndIf
  \EndIf
\EndFunction

\State \textit{\# Typecheck all statements in program}
\For{each statement $s$ in $P$}
  \State \Call{typecheck\_statement}{$s$}
\EndFor

\State \textit{\# Final check: output label matches expected}
\State $L_{\text{result}} \gets$ \Call{join\_all\_labels}{$\Gamma$} \textit{\# Join all variable labels}
\If{$\neg (L_{\text{result}} \leq L_{\text{out}})$}
  \State append("OUTPUT\_LABEL\_MISMATCH: result not $\leq$ expected output" to errors)
\EndIf

\State Output: (errors == []) $\Rightarrow$ PASS, else FAIL(errors)
\end{algorithmic}

\textbf{Role in Design}: 
- Before executing any white-box program, server runs TYPECHECK
- If PASS: program certified secure, execute
- If FAIL: reject program, log error, no execution occurs
- Ensures: no information flows from secret inputs to public outputs without authorization

\textbf{Complexity}: $O(|P|)$ where $|P|$ is program size. Single pass through program statements.

\subsection{Label Join Operation}

\textbf{Purpose}: Compute most restrictive label consistent with multiple labeled data sources. When program reads data from multiple labels, combine labels to get output label. Join enforces: only readers authorized for ALL input labels can access output.

\textbf{Formal Definition}:

For labels $L_1 = \{o_1: R_1\}, L_2 = \{o_2: R_2\}$:

\begin{equation}
L_1 \sqcup L_2 = \{(o_1: R_1 \cap R_2), (o_2: R_1 \cap R_2)\}
\end{equation}

\textit{Reader intersection}: Only principals in $R_1 \cap R_2$ (both reader sets) can read joined result.

\textbf{Pseudocode}:

\begin{algorithmic}
\Function{label\_join}{$L_1$, $L_2$}
  \State result $\gets \{\}$ \textit{\# New label}
  \For{each owner $o$ in $\{o_1, o_2\}$}
    \State $R_1 \gets$ \Call{readers}{$L_1$, $o$}
    \State $R_2 \gets$ \Call{readers}{$L_2$, $o$}
    \State $R \gets R_1 \cap R_2$ \textit{\# Intersection}
    \State add$(o: R)$ to result
  \EndFor
  \State return result
\EndFunction

\Function{label\_join\_multiple}{$L_1, L_2, \ldots, L_k$}
  \State result $\gets L_1$
  \For{$i = 2$ to $k$}
    \State result $\gets$ \Call{label\_join}{result, $L_i$}
  \EndFor
  \State return result
\EndFunction
\end{algorithmic}

\textbf{Example}:
\begin{itemize}
\item $L_1 = \text{patient\_1}: \{\text{doctor\_1}, \text{patient\_1}\}$ (patient data readable by doctor and patient)
\item $L_2 = \text{hospital}: \{\text{hospital}, \text{analyst}\}$ (hospital database readable by hospital and analyst)
\item Join: For each owner, intersect readers:
  \begin{itemize}
  \item patient\_1 readers: $\{\text{doctor\_1}, \text{patient\_1}\} \cap \{\text{hospital}, \text{analyst}\} = \{\}$ (empty)
  \item hospital readers: $\{\text{doctor\_1}, \text{patient\_1}\} \cap \{\text{hospital}, \text{analyst}\} = \{\}$ (empty)
  \end{itemize}
\item Result: $L_{\text{join}} = \text{patient\_1}: \{\}, \text{hospital}: \{\}$ (nobody can read combined output—most restrictive)
\end{itemize}

\textbf{Role in Design}:
- Used in black-box program labeling: output = join of all inputs
- Used in white-box type-checking: combined labels for expressions using multiple inputs
- Ensures: most restrictive label applied; security guaranteed by conservatism

\subsection{Label Ordering Partial Order}

\textbf{Purpose}: Define when one label is more restrictive (can safely flow to) another. Label ordering $L_1 \leq L_2$ encodes: if $L_1$ restricts more, then information with $L_1$ can flow to variables expecting $L_2$.

\textbf{Formal Definition}:

\begin{equation}
L_1 \leq L_2 \text{ iff } \forall o \in \text{principals}, R_1(o) \supseteq R_2(o)
\end{equation}

Interpretation: $L_1$ more restrictive iff readers in $L_1$ are superset of readers in $L_2$ (fewer readers for same principal).

\textbf{Properties}:
\begin{itemize}
\item \textbf{Reflexive}: $L \leq L$ (label ordering itself)
\item \textbf{Transitive}: If $L_1 \leq L_2$ and $L_2 \leq L_3$, then $L_1 \leq L_3$
\item \textbf{Antisymmetric}: If $L_1 \leq L_2$ and $L_2 \leq L_1$, then $L_1 = L_2$
\end{itemize}

Forms a partial order over labels.

\textbf{Pseudocode}:

\begin{algorithmic}
\Function{label\_leq}{$L_1$, $L_2$}
  \For{each owner $o$}
    \State $R_1 \gets$ \Call{readers}{$L_1$, $o$}
    \State $R_2 \gets$ \Call{readers}{$L_2$, $o$}
    \If{$\neg (R_1 \supseteq R_2)$}
      \State return FALSE \textit{\# Readers mismatch; $L_1 \not\leq L_2$}
    \EndIf
  \EndFor
  \State return TRUE \textit{\# All owners satisfy $R_1 \supseteq R_2$}
\EndFunction
\end{algorithmic}

\textbf{Example}:
\begin{itemize}
\item $L_1 = \text{patient}: \{\}$ (no readers; most restrictive)
\item $L_2 = \text{patient}: \{\text{doctor}\}$ (doctor can read)
\item Is $L_1 \leq L_2$? Check: readers in $L_1$ $\supseteq$ readers in $L_2$? $\{\} \supseteq \{\text{doctor}\}$? NO
\item Result: $L_1 \not\leq L_2$ (restrictive $L_1$ cannot flow to less-restrictive $L_2$—counterintuitive but correct!)
\item Instead: $L_2 \leq L_1$? Check: $\{\text{doctor}\} \supseteq \{\}$? YES. So less-restrictive $L_2$ can flow to $L_1$.
\end{itemize}

\textbf{Role in Design}:
- Type-checker uses label ordering for explicit flow check: assignment $x := e$ requires $\text{label}(e) \leq \text{label}(x)$
- Ensures: information cannot flow from less-restrictive sources to more-restrictive sinks
- Counterintuitive direction: less-restrictive $\to$ more-restrictive allowed; more-restrictive $\to$ less-restrictive blocked

\subsection{Hospital Scenario Label Lattice}

\textbf{Purpose}: Document concrete label instantiations for the hospital medical research scenario. Shows how abstract Myers-Liskov framework applies to concrete system with patients, doctors, researchers, analyzers.

\textbf{Principals}:
\begin{itemize}
\item $P_1, P_2, \ldots$ (patients)
\item $D_1, D_2, \ldots$ (treating doctors)
\item $H$ (hospital institution)
\item $R_1, R_2, \ldots$ (researchers)
\item $A_1, A_2, \ldots$ (analyzer AI systems)
\item $\text{Server}$ (hosting server)
\item $\text{Public}$ (external observers)
\end{itemize}

\textbf{Key Labels}:

\begin{tabular}{|l|l|l|}
\hline
\textbf{Data} & \textbf{Label} & \textbf{Interpretation} \\
\hline
Patient data & $P_i: \{D_i, P_i\}$ & Only doctor + patient read \\
\hline
Hospital DB & $H: \{H\}$ & Only hospital reads \\
\hline
Processed data & $A: \{\}$ (after analysis) & Only analyzer sees (most restrictive) \\
\hline
Anonymized cohort & $R: \{R, \text{Public}\}$ & Researchers + public after de-id \\
\hline
Aggregates & $R: \{\text{Public}\}$ & Public accessible after aggregation \\
\hline
\end{tabular}

\textbf{Label Transitions}:

\begin{enumerate}
\item Upload: $L_0 = P_1: \{D_1, P_1\}$ (patient data)
\item Declassify for research: $L_1 = P_1: \{D_1, P_1, A\}$ (grant analyzer access)
\item De-identification: $L_2 = R: \{R, \text{Public}\}$ (ownership transfer to researcher; de-id justifies public readers)
\item Publication: $L_3 = R: \{\text{Public}\}$ (pure public aggregates)
\end{enumerate}

\textbf{Role in Design}:
- Concrete instantiation of abstract framework
- Shows how to model hospital medical scenario
- Demonstrates label transitions at each declassification step
- Provides reference for implementing server

\subsection{Access Control Matrix}

\textbf{Purpose}: Specify which principals can read/write which data based on labels. Access control derived from label definitions, enforced by server READ/WRITE commands.

\textbf{Rules}:

For data with label $L = o: R$ (owner $o$, readers $R$):
\begin{itemize}
\item \textbf{Read Access}: Principal $p$ can read iff $p = o$ OR $p \in R$
\item \textbf{Write Access}: Principal $p$ can write iff $p = o$ (only owner)
\item \textbf{Declassify}: Principal $p$ can declassify iff $p = o$ (only owner)
\end{itemize}

\textbf{Pseudocode}:

\begin{algorithmic}
\Function{can\_read}{principal $p$, label $L$}
  \State $o \gets$ \Call{owner}{$L$}
  \State $R \gets$ \Call{readers}{$L$}
  \State return $(p = o)$ \textbf{or} $(p \in R)$
\EndFunction

\Function{can\_write}{principal $p$, label $L$}
  \State $o \gets$ \Call{owner}{$L$}
  \State return $(p = o)$
\EndFunction

\Function{can\_declassify}{principal $p$, label $L$}
  \State return \Call{can\_write}{$p$, $L$} \textit{\# Same as write}
\EndFunction
\end{algorithmic}

\textbf{Example}:

For data labeled $P_1: \{D_1\}$ (patient owns, doctor reads):
\begin{itemize}
\item Can $P_1$ read? $P_1 = \text{owner}$ ✓ YES
\item Can $D_1$ read? $D_1 \in \text{readers}$ ✓ YES
\item Can $D_2$ read? $D_2 \neq \text{owner}$ and $D_2 \notin \text{readers}$ ✗ NO
\item Can $P_1$ write? $P_1 = \text{owner}$ ✓ YES
\item Can $D_1$ write? $D_1 \neq \text{owner}$ ✗ NO
\item Can $P_1$ declassify? $P_1 = \text{owner}$ ✓ YES
\end{itemize}

\textbf{Role in Design}:
- Implements server authorization logic
- Enforced by READ, WRITE, DECLASSIFY commands
- Ensures only authorized principals access data

\subsection{Audit Trail Schema}

\textbf{Purpose}: Document format for audit logs. Every security-relevant operation (UPLOAD, READ, DECLASSIFY) logged with timestamp, principal, action, parameters, signature. Enables post-hoc breach investigation and accountability.

\textbf{Schema}:

\begin{verbatim}
audit_entry = {
  timestamp: ISO8601_datetime,
  principal: principal_id,
  action: "UPLOAD" | "READ" | "DECLASSIFY" | "ANALYZE_WHITEBOX" | "RUN_BLACKBOX",
  data_id: string,
  old_label: label (if applicable),
  new_label: label (if applicable),
  parameters: {
    (action-specific fields)
  },
  justification: string (if DECLASSIFY),
  result: "SUCCESS" | "FAILURE",
  error_message: string (if FAILURE),
  signature: SHA256(entry_without_signature)
}
\end{verbatim}

\textbf{Examples}:

\textbf{UPLOAD}:
\begin{verbatim}
{
  timestamp: "2026-04-20T15:04:26Z",
  principal: "patient_1",
  action: "UPLOAD",
  data_id: "P1_CT_SCAN_001",
  new_label: "patient_1:{doctor_1, patient_1}",
  result: "SUCCESS"
}
\end{verbatim}

\textbf{DECLASSIFY}:
\begin{verbatim}
{
  timestamp: "2026-04-20T15:05:00Z",
  principal: "patient_1",
  action: "DECLASSIFY",
  data_id: "P1_CT_SCAN_001",
  old_label: "patient_1:{doctor_1, patient_1}",
  new_label: "patient_1:{doctor_1, patient_1, analyzer}",
  justification: "Authorize medical AI analysis for cancer research",
  result: "SUCCESS"
}
\end{verbatim}

\textbf{Role in Design}:
- Complete record of all security operations
- Timestamp enables ordering
- Principal + signature enable accountability (audit entry signed by principal's key)
- Justification required for declassification (business decision documentation)
- Error messages for failed operations help debugging

\subsection{Design Pattern: Label Polymorphism}

\textbf{Purpose}: Show how labels vary at different program points. Single variable can have different labels in different contexts. Enables fine-grained information flow analysis.

\textbf{Pattern}:

In white-box program:
\begin{verbatim}
x = read_file_1();         // x has label L1 (from file_1)
y = read_file_2();         // y has label L2 (from file_2)
z = x + y;                 // z has label L1 ⊔ L2 (join of both)
if (z > threshold) {       // program counter label = L1 ⊔ L2
  result = true;           // result constrained by (L1 ⊔ L2)
}
\end{verbatim}

Type-checking tracks:
\begin{itemize}
\item \textit{Entry}: $x: L_1$, $y: L_2$
\item \textit{After $z = x + y$}: $z: L_1 \sqcup L_2$
\item \textit{Inside if}: $L_{pc} = L_1 \sqcup L_2$
\item \textit{Assignment $\text{result} = \text{true}$}: $L_{pc} \sqcup \text{label}(\text{true}) = L_1 \sqcup L_2 \sqcup \text{public} = L_1 \sqcup L_2 \leq \text{label}(\text{result})$
\end{itemize}

\textbf{Role in Design}:
- Shows how implicit flow control prevents information leaks through conditionals
- Label annotations on variables must be conservative (least restrictive possible)
- Type-checker proves: if annotations satisfy rules, program secure

\end{document}
